\chapter{Mécanismes réactionnels des complexes}\label{ch:b3:complex-mechanisms}

Dissolvons un sel de chrome(III) dans de l'eau marquée à l'oxygène 18~: il
faut des jours pour que l'eau marquée remplace les molécules d'eau liées au
métal~; avec un sel de cuivre(II), l'échange est terminé dans le temps du
mélange. La réaction est la même, une molécule d'eau qui quitte un ion
métallique et une autre qui prend sa place, et pourtant sa vitesse varie
énormément d'un ion à l'autre, et le nombre d'électrons $d$ prévoit une
grande partie de ces variations. La façon dont les ligands sont remplacés
décide de la fabrication des anticancéreux à base de platine et de leur
mode d'action~; la façon dont les électrons passent d'un ion métallique à
l'autre décide de la vitesse de la respiration et de toutes les batteries.
Ce chapitre traite des deux familles de réactions des complexes~: la
substitution et le transfert d'électron, ce dernier avec la théorie de
Rudolph Marcus.

\begin{recall}
Le volume de deuxième année a décrit les complexes, la denticité des
ligands, les ligands pontants, les isomères $fac$ et $mer$, l'échange de
ligands comme acte élémentaire, et les énergies de stabilisation du champ
cristallin avec les configurations haut spin et bas spin~; le volume de
première année, les mécanismes avec les approximations de l'état quasi
stationnaire et de l'étape cinétiquement déterminante. Le
\cref{ch:b3:rate-theories} a donné l'\omterm{thm:b3:rate-theories:eyring}{équation d'Eyring} et le sens de
$\Delta^{\ddagger}S^\circ$, le \cref{ch:b3:electrode-kinetics} le transfert
d'électron à une électrode, et le \cref{ch:b3:complex-spectra-magnetism} les
\omterm{def:b3:complex-spectra-magnetism:ligand-field-term}{termes du champ des ligands} et l'effet Jahn--Teller.
\end{recall}

\section{Labilité et inertie}

\begin{definition}[Complexes labiles et inertes]\label{def:b3:complex-mechanisms:labile}
Un complexe est \emph{labile}\index{complexe labile} si ses ligands sont
remplacés rapidement, et c'est un \emph{complexe cinétiquement inerte}\index{complexe cinetiquement inerte@complexe cinétiquement inerte}
s'ils sont remplacés lentement~; suivant Taube, on dit qu'un complexe est
labile lorsque ses réactions avec un ligand à \qty{0.1}{mol/L} et
\qty{25}{\celsius} sont terminées en une minute environ. La labilité est
une propriété cinétique, sans rapport avec la stabilité thermodynamique du
complexe.
\end{definition}

L'échange d'eau, $\mathrm{[M(H_2O)_6]^{n+} + H_2O^{\ast} \to [M(H_2O)_5(H_2O^{\ast})]^{n+} + H_2O}$, est la substitution la plus
simple et fournit une échelle. Les ions des groupes principaux échangent
d'autant plus vite qu'ils sont plus gros et moins chargés. Parmi les ions
des métaux de transition, ceux de configuration $d^3$ (\ce{Cr^{3+}}) ou $d^6$
bas spin (\ce{Co^{3+}}, \ce{Rh^{3+}}, \ce{Ir^{3+}}) sont inertes, comme la
plupart des métaux des deuxième et troisième lignes, tandis que $d^9$
(\ce{Cu^{2+}}) et $d^4$ haut spin (\ce{Cr^{2+}}), déformés par l'effet
Jahn--Teller, sont extrêmement \omterm{def:b3:complex-mechanisms:labile}{labiles}.

\begin{proposition}[Énergie d'activation de champ des ligands]\label{prop:b3:complex-mechanisms:lfae}
Si une substitution passe par un état de transition pentacoordiné en
pyramide à base carrée, la perte de stabilisation du champ des ligands
subie pour l'atteindre (l'énergie d'activation de champ des ligands) est,
dans le modèle le plus simple, purement $\sigma$, maximale pour $d^6$ bas
spin ($4\,Dq$), puis pour $d^3$ et $d^8$ ($2\,Dq$), nulle pour $d^0$, $d^5$
haut spin et $d^{10}$, et négative pour $d^4$ haut spin et $d^9$.
\end{proposition}

\begin{proof}
On modélise chaque ligand comme élevant une orbitale $d$ de $e_\sigma$ fois
le carré de l'amplitude angulaire de l'orbitale dans sa direction,
normalisée à 1 sur l'axe du lobe~: un ligand axial élève $d_{z^2}$ de 1, un
ligand équatorial élève $d_{z^2}$ de $\frac14$ et $d_{x^2-y^2}$ de $\frac34$~; les orbitales $t_{2g}$
pointent entre les ligands et ne reçoivent rien. L'octaèdre élève $d_{z^2}$
et $d_{x^2-y^2}$ de $3e_\sigma$ chacune ($\Delta_{\mathrm o} = 3e_\sigma = 10Dq$)~; la pyramide à base
carrée, privée d'un ligand axial, élève $d_{z^2}$ de $2e_\sigma$ et $d_{x^2-y^2}$
de $3e_\sigma$. La stabilisation de $n$ électrons est leur énergie moins
$n/5$ fois la somme des cinq niveaux (le barycentre). Pour $d^6$ bas spin
($t_{2g}^6$)~: $0 - \frac65 \times 6e_\sigma$ dans l'octaèdre, $0 - \frac65 \times 5e_\sigma$ dans la
pyramide~: une perte de $1.2e_\sigma = 4Dq$. Pour $d^3$~: $\frac35e_\sigma =
2Dq$~; pour $d^8$ ($t_{2g}^6e_g^2$)~: $(5 - \frac85 \times 5) - (6 - \frac85 \times 6) = 0.6e_\sigma$. Pour $d^9$~: $(7 - 9) - (9 -
10.8) = -0.2e_\sigma$.
\end{proof}

\begin{omfigure}
\begin{tikzpicture}
\begin{axis}[width=0.72\linewidth, height=5.4cm, xmin=-0.6, xmax=10.6, ymin=-1.2, ymax=4.6,
  axis x line=bottom, axis y line=left, ycomb, xtick={0,1,2,3,4,5,6,7,8,9,10},
  xlabel={$n$ ($d^n$)}, ylabel={énergie d'activation / $Dq$}, tick label style={font=\small},
  label style={font=\small}]
\draw[black!40] (axis cs:-0.6,0) -- (axis cs:10.6,0);
\addplot[omDef, line width=4pt, xshift=-2.5pt] table[x=n, y=hs] {figdata/out/bachelor-3-ligand-field-activation.dat};
\addplot[omThm, line width=4pt, xshift=2.5pt, unbounded coords=discard] table[x=n, y=ls] {figdata/out/bachelor-3-ligand-field-activation.dat};
\end{axis}
\end{tikzpicture}

{\small Énergies d'activation de champ des ligands pour un chemin
dissociatif passant par une pyramide à base carrée, dans le modèle purement
$\sigma$ de la proposition (en bleu~: haut spin~; en rouge~: bas spin). Les
plus grandes valeurs, pour $d^6$ bas spin, $d^3$ et $d^8$, correspondent aux
ions inertes ou lents (\ce{Co^{3+}}, \ce{Cr^{3+}}, \ce{Ni^{2+}})~; les
valeurs négatives, pour $d^4$ et $d^9$, aux ions Jahn--Teller très \omterm{def:b3:complex-mechanisms:labile}{labiles}.}
\end{omfigure}

\begin{definition}[Volume d'activation]\label{def:b3:complex-mechanisms:activation-volume}
Le \emph{volume d'activation}\index{volume d'activation} $\Delta V^{\ddagger}$
d'une étape élémentaire est la différence entre le volume molaire partiel
du \omterm{def:b3:rate-theories:activated-complex}{complexe activé} et ceux des réactifs.
\end{definition}

\begin{proposition}[Influence de la pression sur une constante de vitesse]\label{prop:b3:complex-mechanisms:pressure}
À température constante, $\displaystyle\Bigl(\frac{\partial\ln k}{\partial p}\Bigr)_T = -\frac{\Delta V^{\ddagger}}{RT}$.
\end{proposition}

\begin{proof}
D'après l'\omterm{thm:b3:rate-theories:eyring}{équation d'Eyring}, $\ln k = \text{cste} - \Delta^{\ddagger}G/RT$, et $(\partial\Delta^{\ddagger}G/\partial p)_T = \Delta V^{\ddagger}$, comme pour
toute enthalpie libre.
\end{proof}

Une étape au cours de laquelle un ligand part a $\Delta V^{\ddagger} > 0$~;
une étape au cours de laquelle un ligand entre a $\Delta V^{\ddagger}
< 0$. Mesuré à quelques centaines de mégapascals, le \omterm{def:b3:complex-mechanisms:activation-volume}{volume d'activation}
est le diagnostic le plus direct d'un mécanisme de substitution, moins
brouillé par la solvatation que l'\omterm{def:b3:rate-theories:activation}{entropie d'activation}.

\section{Mécanismes de substitution}

\begin{definition}[Mécanismes dissociatif, associatif et d'échange]\label{def:b3:complex-mechanisms:mechanisms}
Dans un \emph{mécanisme dissociatif}\index{mecanisme dissociatif@mécanisme dissociatif} (D), le ligand
partant s'en va d'abord, ce qui donne un intermédiaire de coordinence plus
faible~; dans un \emph{mécanisme associatif}\index{mecanisme associatif@mécanisme associatif} (A), le ligand
entrant se lie d'abord, ce qui donne un intermédiaire de coordinence plus
élevée~; dans un \emph{mécanisme d'échange}\index{mecanisme d'echange@mécanisme d'échange} (I), les deux
ligands s'échangent en une seule étape, à dominante de rupture de liaison
($I_d$) ou de formation de liaison ($I_a$).
\end{definition}

\begin{proposition}[Lois de vitesse des mécanismes D et A]\label{prop:b3:complex-mechanisms:d-a-laws}
Pour $\mathrm{ML_5X + Y \to ML_5Y + X}$~: un mécanisme D ($\mathrm{ML_5X} \rightleftharpoons \mathrm{ML_5} + \mathrm X$, $k_1$, $k_{-1}$~; $\mathrm{ML_5} + \mathrm Y
\to \mathrm{ML_5Y}$, $k_2$) donne
\[
  v = \frac{k_1k_2[\mathrm{ML_5X}][\mathrm Y]}{k_{-1}[\mathrm X] + k_2[\mathrm Y]},
\]
du premier ordre par rapport au complexe et indépendante de $[\mathrm Y]$
aux fortes valeurs de $[\mathrm Y]$~; un mécanisme A passant par un
intermédiaire heptacoordiné en régime stationnaire donne $v =
k[\mathrm{ML_5X}][\mathrm Y]$, du premier ordre par rapport à chacun.
\end{proposition}

\begin{proof}
D~: l'approximation de l'état quasi stationnaire pour $\mathrm{ML_5}$ donne $k_1[\mathrm{ML_5X}] = (k_{-1}[\mathrm X] + k_2[\mathrm Y])[\mathrm{ML_5}]$, et $v = k_2[\mathrm{ML_5}][\mathrm Y]$.
Pour $[\mathrm Y]$ élevée, $v \to k_1[\mathrm{ML_5X}]$. A~: $\mathrm{ML_5X} + \mathrm Y \rightleftharpoons \mathrm{ML_5XY}$ ($k_a$, $k_{-a}$), $\mathrm{ML_5XY} \to
\mathrm{ML_5Y} + \mathrm X$ ($k_b$)~; l'état quasi stationnaire donne $v = \frac{k_ak_b}{k_{-a} + k_b}[\mathrm{ML_5X}][\mathrm Y]$.
\end{proof}

Dans l'eau, le ligand entrant est en général présent à une concentration
bien plus faible que le solvant, et la véritable première étape est la
rencontre du complexe et du ligand.

\begin{definition}[Mécanisme d'Eigen--Wilkins]\label{def:b3:complex-mechanisms:eigen-wilkins}
Le \emph{mécanisme d'Eigen--Wilkins}\index{mecanisme d'Eigen--Wilkins@mécanisme d'Eigen--Wilkins} de
substitution d'un aquacomplexe est un pré-équilibre rapide qui forme un
\emph{complexe de sphère externe}\index{complexe de sphere externe@complexe de sphère externe}, dans lequel le ligand
entrant se place contre le complexe, hors de sa sphère de coordination,
suivi de l'échange cinétiquement déterminant d'une molécule d'eau contre ce
ligand.
\end{definition}

\begin{theorem}[Loi de vitesse d'Eigen--Wilkins]\label{thm:b3:complex-mechanisms:eigen-wilkins}
Avec la constante d'équilibre de sphère externe $K_{\mathrm{os}}$ et la
constante de vitesse d'échange $k_i$, la constante de vitesse apparente de
pseudo-premier ordre, pour $[\mathrm Y]$ en excès, est
\[
  k_{\mathrm{obs}} = \frac{K_{\mathrm{os}}k_i[\mathrm Y]}{1 + K_{\mathrm{os}}[\mathrm Y]} .
\]
\end{theorem}

\begin{proof}
Le complexe se répartit entre forme libre et forme de sphère externe, $[\mathrm{M\cdot Y}] = K_{\mathrm{os}}[\mathrm M][\mathrm Y]$,
avec au total $[\mathrm M]_{\mathrm{tot}} = [\mathrm M](1 + K_{\mathrm{os}}[\mathrm Y])$~; la vitesse vaut $k_i[\mathrm{M\cdot Y}] = k_iK_{\mathrm{os}}[\mathrm Y][\mathrm M]_{\mathrm{tot}}/(1 + K_{\mathrm{os}}[\mathrm Y])$.
\end{proof}

Pour $[\mathrm Y]$ faible, la réaction est du second ordre avec $k = K_{\mathrm{os}}k_i$~; $K_{\mathrm{os}}$,
estimée à partir des charges et de la distance de plus courte approche, est
presque la même pour tous les ligands de même charge, si bien que $k_i$,
l'étape d'échange d'eau, contrôle la vitesse~: la substitution sur un
aqua-ion est presque indépendante du ligand entrant, signature d'un
échange dissociatif.

\begin{definition}[Mécanisme de la base conjuguée]\label{def:b3:complex-mechanisms:conjugate-base}
Le \emph{mécanisme de la base conjuguée}\index{mecanisme de la base conjuguee@mécanisme de la base conjuguée} de
l'hydrolyse basique d'un complexe ammine, $\mathrm{[Co(NH_3)_5X]^{2+} + OH^- \to [Co(NH_3)_5(OH)]^{2+} + X^-}$, est
la déprotonation d'un ligand ammine dans un pré-équilibre rapide, suivie de
la perte dissociative de $\mathrm X^-$ par la base conjuguée amido, dont le
ligand $\mathrm{NH_2^-}$ labilise le complexe.
\end{definition}

\begin{proposition}[Loi de vitesse de l'hydrolyse basique]\label{prop:b3:complex-mechanisms:conjugate-base}
Avec la constante de déprotonation $K$ et la constante de vitesse de
dissociation $k$ de la base conjuguée, $v = \dfrac{kK[\text{complexe}]_{\mathrm{tot}}[\ce{OH-}]}{1 + K[\ce{OH-}]}$, du second ordre pour $[\ce{OH-}]$ faible.
\end{proposition}

\begin{proof}
Comme pour la loi d'Eigen--Wilkins~: le complexe se répartit entre ses
formes acide et base conjuguée dans le rapport $1 : K[\ce{OH-}]$, et seule
la seconde réagit.
\end{proof}

La loi de vitesse seule ne distingue pas ce mécanisme d'une attaque directe
de \ce{OH-}~; les arguments sont que l'hydrolyse basique exige un proton
N--H (les complexes qui n'en ont pas ne la présentent pas) et que les
protons des ligands ammine s'échangent avec le solvant bien plus vite que
l'hydrolyse.

\begin{method}[Diagnostiquer un mécanisme de substitution]\label{met:b3:complex-mechanisms:diagnose}
\begin{enumerate}
  \item Dépendance vis-à-vis du groupe entrant~: vitesse indépendante de Y
        (ou saturante), \omterm{def:b3:complex-mechanisms:mechanisms}{mécanisme dissociatif}~; vitesse proportionnelle à
        $[\mathrm Y]$ et sensible à la nature de Y, \omterm{def:b3:complex-mechanisms:mechanisms}{mécanisme associatif}.
  \item \omterm{def:b3:complex-mechanisms:activation-volume}{Volume d'activation}~: positif, dissociatif~; négatif, associatif.
  \item \omterm{def:b3:rate-theories:activation}{Entropie d'activation}~: positive pour D, négative pour A (avec
        prudence, car la solvatation y contribue).
  \item Dépendance vis-à-vis du groupe partant~: une vitesse qui suit la
        force de la liaison M--X indique une rupture de liaison dans l'état
        de transition.
\end{enumerate}
\end{method}

\begin{omfigure}
\begin{tikzpicture}
\begin{axis}[name=pd, omprofile, width=0.46\linewidth, height=4.6cm, xmin=0, xmax=10, ymin=0, ymax=10,
  xlabel={coordonnée de réaction}, ylabel={énergie}, title={\small D~: $\mathrm{ML_5X} \to \mathrm{ML_5} + \mathrm X \to \mathrm{ML_5Y}$}]
\addplot[omDef, very thick, smooth] coordinates {(0,2) (1.5,2.1) (3,7.5) (4.2,5.6) (5.0,5.4) (5.8,5.6) (7,7.0) (8.5,1.6) (10,1.5)};
\node[font=\scriptsize, anchor=north, align=center] at (axis cs:5,5.2) {ML$_5$ (coordinence 5)};
\end{axis}
\begin{axis}[at={(pd.east)}, anchor=west, xshift=1.0cm, omprofile, width=0.46\linewidth, height=4.6cm,
  xmin=0, xmax=10, ymin=0, ymax=10, xlabel={coordonnée de réaction}, ylabel={énergie},
  title={\small A~: $\mathrm{ML_5X} + \mathrm Y \to \mathrm{ML_5XY} \to \mathrm{ML_5Y}$}]
\addplot[omThm, very thick, smooth] coordinates {(0,2) (1.5,2.1) (3,6.5) (4.2,4.4) (5.0,4.2) (5.8,4.4) (7,7.2) (8.5,1.6) (10,1.5)};
\node[font=\scriptsize, anchor=north, align=center] at (axis cs:5,4.0) {ML$_5$XY (coordinence 7)};
\end{axis}
\end{tikzpicture}

{\small Profils énergétiques d'une substitution dissociative et d'une
substitution associative (schématiques)~: toutes deux passent par un
intermédiaire, de coordinence plus faible en D et plus élevée en A. Dans les
\omterm{def:b3:complex-mechanisms:mechanisms}{mécanismes d'échange}, le puits disparaît et il ne reste qu'un seul état de
transition.}
\end{omfigure}

Un complexe octaédrique qui se substitue en passant par une pyramide à base
carrée conserve en général sa configuration ($cis$ reste $cis$)~; un
complexe qui passe par une bipyramide trigonale, comme beaucoup de
complexes du cobalt(III) en hydrolyse basique, peut donner un mélange de
produits $cis$ et $trans$.

\section{Substitution plan carré et effet trans}

Les complexes plans carrés des ions $d^8$ (\ce{Pt^{2+}}, \ce{Pd^{2+}}, \ce{Au^{3+}}, \ce{Ni^{2+}} avec des ligands à champ fort)
sont coordinativement insaturés~: un ligand entrant peut approcher selon la
direction axiale libre. Leurs substitutions sont associatives et passent par
un état de transition pentacoordiné bipyramidal trigonal.

\begin{proposition}[Loi de vitesse à deux termes]\label{prop:b3:complex-mechanisms:square-planar}
La substitution $\mathrm{[PtL_3X] + Y \to [PtL_3Y] + X}$ dans un solvant coordinant S suit
$k_{\mathrm{obs}} = k_1 + k_2[\mathrm Y]$ pour Y en excès.
\end{proposition}

\begin{proof}
Deux chemins associatifs parallèles~: l'attaque directe de Y ($k_2[\mathrm Y]$), et
l'attaque du solvant, présent en large excès (pseudo-premier ordre, $k_1$),
qui donne $[\mathrm{PtL_3S}]$, ensuite rapidement converti par Y. Des chemins
parallèles du premier ordre s'additionnent.
\end{proof}

\begin{definition}[Effet trans, influence trans]\label{def:b3:complex-mechanisms:trans-effect}
L'\emph{effet trans}\index{effet trans} est l'accélération de la
substitution d'un ligand par le ligand situé en trans~: un effet cinétique,
sur l'état de transition. L'\emph{influence trans}\index{influence trans}
est l'affaiblissement, dans l'état fondamental, de la liaison située en
trans d'un ligand, visible sur les longueurs de liaison et les fréquences
de vibration~: un effet thermodynamique.
\end{definition}

L'\omterm{def:b3:complex-mechanisms:trans-effect}{effet trans} suit l'ordre $\ce{CO}$, $\ce{CN-}$, $\ce{C2H4} > \ce{PR3}$, $\ce{H-} > \ce{CH3-} > \ce{I-}$, $\ce{SCN-} >
\ce{Br-} > \ce{Cl-} > \ce{NH3}$, $\ce{H2O}$. Les donneurs $\sigma$ forts
affaiblissent la liaison en trans (\omterm{def:b3:complex-mechanisms:trans-effect}{influence trans})~; les accepteurs $\pi$
forts stabilisent l'état de transition pentacoordiné en prenant de la
\omterm{def:b3:computational-chemistry:dft}{densité électronique} au métal. Les deux augmentent la vitesse.

\begin{method}[Planifier la synthèse d'isomères du platine(II)]\label{met:b3:complex-mechanisms:pt-synthesis}
\begin{enumerate}
  \item À chaque étape, le ligand remplacé est celui qui est en trans du
        ligand d'\omterm{def:b3:complex-mechanisms:trans-effect}{effet trans} le plus fort présent.
  \item Ordonner les additions de sorte que cette règle conduise à
        l'isomère voulu.
\end{enumerate}
\end{method}

\begin{omfigure}
\setchemfig{atom sep=2.2em}
\schemestart
\chemfig{Cl-Pt(-[2]Cl)(-[6]Cl)-Cl}
\arrow{->[\ce{NH3}]}
\chemfig{Cl-Pt(-[2]Cl)(-[6]Cl)-NH_3}
\arrow{->[\ce{NH3}]}
\chemfig{Cl-Pt(-[2]NH_3)(-[6]Cl)-NH_3}
\schemestop

\bigskip
\schemestart
\chemfig{H_3N-Pt(-[2]NH_3)(-[6]NH_3)-NH_3}
\arrow{->[\ce{Cl-}]}
\chemfig{H_3N-Pt(-[2]NH_3)(-[6]NH_3)-Cl}
\arrow{->[\ce{Cl-}]}
\chemfig{Cl-Pt(-[2]NH_3)(-[6]NH_3)-Cl}
\schemestop

{\small L'\omterm{def:b3:complex-mechanisms:trans-effect}{effet trans} à l'œuvre (charges omises). En haut~: à partir de
$\ce{[PtCl4]^{2-}}$, la deuxième molécule d'ammoniac remplace un chlorure en
trans d'un chlorure (Cl$^-$ a l'\omterm{def:b3:complex-mechanisms:trans-effect}{effet trans} le plus fort), ce qui donne
l'isomère $cis$, le cisplatine. En bas~: à partir de $\ce{[Pt(NH3)4]^{2+}}$,
le deuxième chlorure remplace l'ammoniac en trans du premier chlorure, ce
qui donne l'isomère $trans$.}
\end{omfigure}

\section{Transfert d'électron}

\begin{definition}[Transfert d'électron de sphère externe et de sphère interne]\label{def:b3:complex-mechanisms:electron-transfer}
Dans un \emph{transfert d'électron de sphère externe}\index{transfert d'electron de sphere externe@transfert d'électron de sphère externe},
l'électron passe entre deux complexes dont les sphères de coordination
restent intactes~; dans un \emph{transfert d'électron de sphère interne}\index{transfert d'electron de sphere interne@transfert d'électron de sphère interne},
il passe par un ligand qui ponte les deux métaux dans un complexe
binucléaire transitoire. Une \emph{réaction d'auto-échange}\index{reaction d'auto-echange@réaction d'auto-échange} est un
transfert d'électron entre les deux degrés d'oxydation d'un même couple,
sans changement chimique net, comme $\ce{[Fe(H2O)6]^{3+} + [Fe(H2O)6]^{2+}}$.
\end{definition}

Henry Taube a mis en évidence le chemin de sphère interne en 1953 grâce à
une réaction choisie pour que le produit raconte son histoire. Le
cobalt(III) est inerte, le cobalt(II) \omterm{def:b3:complex-mechanisms:labile}{labile}~; le chrome(II) est \omterm{def:b3:complex-mechanisms:labile}{labile}, le
chrome(III) inerte. Quand $\ce{[Co(NH3)5Cl]^{2+}}$ est réduit par $\ce{[Cr(H2O)6]^{2+}}$ en
milieu acide, tout le chrome(III) formé porte le chlorure, sous forme de $\ce{[Cr(H2O)5Cl]^{2+}}$,
même en présence de chlorure libre~: le chlorure a forcément ponté les deux
métaux, s'est retrouvé lié au chrome(III) inerte dès le passage de
l'électron, et a été libéré par le cobalt(II) \omterm{def:b3:complex-mechanisms:labile}{labile}.

\begin{omfigure}
\begin{tikzpicture}[font=\small]
  \begin{scope}[scale=0.8]
    \pic[transform shape] (a) at (0,0) {omoctahedron};
    \pic[transform shape] (b) at (2.0,0) {omoctahedron};
  \end{scope}
  \node[anchor=north east] at (0.3,-1.0) {$\mathrm{Co^{III}(NH_3)_5}$};
  \node[anchor=north west] at (1.3,-1.0) {$\mathrm{Cr^{II}(H_2O)_5}$};
  \node[omThm, anchor=south] at (0.8,0.15) {Cl};
  \draw[->, thick, omDef] (1.6,1.75) to[bend right=40] node[midway, above] {e$^-$} (0,1.75);
\end{tikzpicture}

{\small L'intermédiaire ponté de Taube (schéma)~: le chlorure, encore lié au
cobalt(III), entre dans la sphère de coordination du chrome(II) en
remplaçant une molécule d'eau~; l'électron traverse le pont, et le chlorure
reste sur le chrome(III), désormais inerte.}
\end{omfigure}

Un électron saute en environ $10^{-16}$ s, bien plus vite que les noyaux ne
se déplacent~: d'après le \omterm{thm:b3:electronic-spectroscopy:franck-condon}{principe de Franck--Condon}
(\cref{ch:b3:electronic-spectroscopy}), le transfert ne peut se produire que
dans une configuration nucléaire où réactifs et produits ont la même
énergie. Pour l'auto-échange du fer, les liaisons \ce{Fe^{III}}--O sont plus
courtes que les liaisons \ce{Fe^{II}}--O~; avant que l'électron puisse
passer, les deux complexes doivent se déformer jusqu'à une géométrie
intermédiaire commune, ce qui coûte de l'énergie.

\section{Théorie de Marcus}

\begin{definition}[Énergie de réorganisation]\label{def:b3:complex-mechanisms:reorganisation}
L'\emph{énergie de réorganisation}\index{energie de reorganisation@énergie de réorganisation} $\lambda$
d'un transfert d'électron est l'énergie nécessaire pour amener les
réactifs, sans transférer l'électron, à la configuration nucléaire
(longueurs de liaison des complexes et orientations du solvant environnant)
des produits. C'est la somme d'une contribution interne, due aux liaisons,
et d'une contribution externe, due au solvant.
\end{definition}

\begin{theorem}[Équation de Marcus]\label{thm:b3:complex-mechanisms:marcus}
Si les enthalpies libres des réactifs et des produits sont des paraboles de
même courbure en fonction d'une coordonnée de réaction $x$ (0 à l'équilibre
des réactifs, 1 à celui des produits), l'\omterm{def:b3:rate-theories:activation}{enthalpie libre d'activation} d'un
transfert d'électron d'enthalpie libre standard de réaction $\Delta G^\circ$
et d'\omterm{def:b3:complex-mechanisms:reorganisation}{énergie de réorganisation} $\lambda$ vaut
\[
  \Delta G^{\ddagger} = \frac{(\lambda + \Delta G^\circ)^2}{4\lambda} .
\]
C'est l'\emph{équation de Marcus}\index{equation de Marcus@équation de Marcus}.
\end{theorem}

\begin{proof}
Posons $G_R = \lambda x^2$ et $G_P = \lambda(x - 1)^2 + \Delta G^\circ$~: la courbure est fixée par $G_R(1) = \lambda$,
la définition de $\lambda$. Le transfert a lieu là où les courbes se
croisent~: $\lambda x^2 = \lambda x^2 - 2\lambda x + \lambda +
\Delta G^\circ$, donc $x^\ast = (\lambda + \Delta G^\circ)/2\lambda$, et $\Delta G^{\ddagger} = G_R(x^\ast) = (\lambda + \Delta G^\circ)^2/4\lambda$.
\end{proof}

\begin{corollary}[Région inversée]\label{cor:b3:complex-mechanisms:inverted}
À $\lambda$ fixé, la vitesse augmente à mesure que la réaction devient plus
exergonique, jusqu'à $-\Delta G^\circ = \lambda$, où $\Delta G^{\ddagger} = 0$~; au-delà, elle
diminue de nouveau.
\end{corollary}

\begin{proof}
$\Delta G^{\ddagger}$ est une parabole en $\Delta G^\circ$ dont le minimum, nul, est atteint en $\Delta G^\circ = -\lambda$~;
pour $\Delta G^\circ < -\lambda$, elle croît de nouveau.
\end{proof}

\begin{definition}[Région inversée]\label{def:b3:complex-mechanisms:inverted}
La \emph{région inversée}\index{region inversee@région inversée} du transfert
d'électron est le domaine de forces motrices $-\Delta G^\circ > \lambda$ dans lequel la vitesse
diminue à mesure que la réaction devient plus exergonique.
\end{definition}

\begin{omfigure}
\begin{tikzpicture}
\begin{axis}[name=mp, width=0.47\linewidth, height=5.6cm, xmin=-0.6, xmax=2.2, ymin=-160, ymax=120,
  axis lines=left, xlabel={coordonnée de réaction $x$}, ylabel={$G$ / \unit{kJ.mol^{-1}}},
  tick label style={font=\small}, label style={font=\small}, xtick={0,1,2}]
\addplot[black, thick] table[x=x, y=GR] {figdata/out/bachelor-3-marcus-parabolas.dat};
\addplot[omDef!50, thick] table[x=x, y=P0] {figdata/out/bachelor-3-marcus-parabolas.dat};
\addplot[omDef, thick] table[x=x, y=P50] {figdata/out/bachelor-3-marcus-parabolas.dat};
\addplot[omThm, thick] table[x=x, y=P100] {figdata/out/bachelor-3-marcus-parabolas.dat};
\addplot[omThm!60, thick, dashed] table[x=x, y=P150] {figdata/out/bachelor-3-marcus-parabolas.dat};
\node[font=\scriptsize, anchor=north] at (axis cs:0,-2) {R};
\end{axis}
\begin{axis}[at={(mp.east)}, anchor=west, xshift=1.4cm, width=0.47\linewidth, height=5.6cm,
  xmin=0, xmax=250, ymin=0, ymax=11.5, axis lines=left, xlabel={$-\Delta G^\circ$ / \unit{kJ.mol^{-1}}},
  ylabel={$\log(k/\unit{s^{-1}})$}, tick label style={font=\small}, label style={font=\small}]
\addplot[omDef, very thick] table[x=mdG, y=logk] {figdata/out/bachelor-3-marcus-rate.dat};
\draw[black!40, dashed] (axis cs:100,0) -- (axis cs:100,11);
\node[font=\scriptsize, anchor=south east] at (axis cs:95,1) {normale};
\node[font=\scriptsize, anchor=south west] at (axis cs:105,1) {inversée};
\end{axis}
\end{tikzpicture}

{\small Théorie de Marcus avec $\lambda = \qty{100}{kJ/mol}$ (modèle). À
gauche~: la parabole des réactifs (en noir) et les paraboles des produits
pour $\Delta G^\circ = 0$ et $-50$ (en bleu), $-100 = -\lambda$ (en rouge) et
\qty{-150}{kJ/mol} (en tirets, \omterm{def:b3:complex-mechanisms:inverted}{région inversée})~; leur point
d'intersection, l'état de transition, descend jusqu'au bas de la courbe des
réactifs pour $-\Delta G^\circ = \lambda$ et remonte au-delà. À droite~: le
logarithme de la constante de vitesse (facteur préexponentiel du modèle
\qty{e11}{s^{-1}}) passe par un maximum en $-\Delta G^\circ = \lambda$.}
\end{omfigure}

\begin{proposition}[Énergie de réorganisation de sphère externe]\label{prop:b3:complex-mechanisms:outer-reorganisation}
Pour deux réactifs sphériques de rayons $a_1$, $a_2$ à une distance $d$ dans
un solvant d'indice de réfraction $n$ et de permittivité relative
$\varepsilon_s$, la contribution du solvant à $\lambda$ est proportionnelle à $\bigl(\frac{1}{2a_1} + \frac{1}{2a_2} - \frac1d\bigr)\bigl(\frac{1}{n^2} - \frac{1}{\varepsilon_s}\bigr)$.
\end{proposition}

\admitted

L'eau, très polaire ($\varepsilon_s$ grande) mais d'indice de réfraction
ordinaire, donne une grande réorganisation du solvant~; les solvants
apolaires, une petite. Les gros réactifs se réorganisent moins~: c'est
pourquoi les protéines de transfert d'électron enfouissent leurs sites
métalliques à l'intérieur de la protéine.

\begin{theorem}[Relation croisée de Marcus]\label{thm:b3:complex-mechanisms:cross-relation}
Pour une réaction croisée $\mathrm{A_{ox} + B_{red} \to A_{red} + B_{ox}}$ de constante d'équilibre $K_{12}$,
entre des couples dont les constantes de vitesse d'auto-échange sont
$k_{11}$ et $k_{22}$,
\[
  k_{12} = \sqrt{k_{11}k_{22}K_{12}f}, \qquad \ln f = \frac{(\ln K_{12})^2}{4\ln(k_{11}k_{22}/Z^2)},
\]
où $Z$ est le facteur de fréquence de collision~; $f \approx 1$ lorsque
$K_{12}$ n'est pas trop grande. C'est la \emph{relation croisée de Marcus}\index{relation croisee de Marcus@relation croisée de Marcus}.
\end{theorem}

\begin{proof}[Démonstration partielle]
Supposons que l'\omterm{def:b3:complex-mechanisms:reorganisation}{énergie de réorganisation} de la réaction croisée soit la
moyenne de celles des auto-échanges, $\lambda_{12} = (\lambda_{11} + \lambda_{22})/2$ (chaque réactif apporte la
moitié de la réorganisation de son propre auto-échange), et que toutes les
constantes de vitesse s'écrivent $k = Z\eu^{-\Delta G^{\ddagger}/RT}$. Alors
$\Delta G_{12}^{\ddagger} = \frac{\lambda_{12}}{4}\bigl(1 + \frac{\Delta G^\circ_{12}}{\lambda_{12}}\bigr)^2 = \frac{\lambda_{12}}{4} + \frac{\Delta G^\circ_{12}}{2} + \frac{(\Delta G^\circ_{12})^2}{4\lambda_{12}}$. Le premier terme donne
$\sqrt{k_{11}k_{22}}$ (puisque $\Delta G^{\ddagger}_{ii} = \lambda_{ii}/4$), le deuxième $\sqrt{K_{12}}$ (avec $\Delta G^\circ_{12} = -RT\ln K_{12}$), le
troisième le facteur $f$.
\end{proof}

\begin{method}[Vitesse d'une réaction croisée à partir de données d'auto-échange]\label{met:b3:complex-mechanisms:marcus-estimate}
\begin{enumerate}
  \item Trouver les constantes d'auto-échange $k_{11}$, $k_{22}$ des deux couples.
  \item Calculer $K_{12}$ à partir des potentiels standard~: $\ln K_{12} = nF(E^\circ_{\text{oxydant}} - E^\circ_{\text{réducteur}})/RT$.
  \item $k_{12} = \sqrt{k_{11}k_{22}K_{12}}$ (avec $f$ si $K_{12}$ est très grande). Un accord à un
        petit facteur près avec la valeur mesurée indique un mécanisme de
        sphère externe.
\end{enumerate}
\end{method}

Marcus a publié sa théorie en 1956~; la \omterm{def:b3:complex-mechanisms:inverted}{région inversée}, d'abord mise en
doute, a été observée en 1984 par Gerhard Closs et John Miller dans des
molécules où le donneur et l'accepteur sont maintenus à distance fixe par
un espaceur rigide, si bien que la diffusion ne peut pas la masquer.

\begin{inthelab}[L'échange d'eau suivi par RMN de l'oxygène 17]
La résonance de l'oxygène 17 de l'eau liée à un ion paramagnétique et celle
de l'eau libre sont élargies par l'échange entre elles. L'enregistrement de
la largeur de raie du signal de l'eau libre en fonction de la température,
dans une solution du perchlorate du métal enrichie en \ce{^{17}O}, donne la
constante de vitesse d'échange et, avec une sonde haute pression, son
\omterm{def:b3:complex-mechanisms:activation-volume}{volume d'activation}.
\end{inthelab}

\begin{safety}{\ghs[0.8cm]{GHS05}\ghs[0.8cm]{GHS06}\ghs[0.8cm]{GHS08}}
% ledger: ghs4:K2PtCl4
Tétrachloroplatinate(II) de potassium~: toxique en cas d'ingestion,
provoque des lésions oculaires graves, sensibilisant respiratoire et
cutané. Les complexes du platine se manipulent avec des gants sous une
hotte~; le cisplatine et ses analogues sont des médicaments cytotoxiques et
sont manipulés comme tels.
\end{safety}

\begin{history}[Taube et Marcus]
Henry Taube a classé les complexes en \omterm{def:b3:complex-mechanisms:labile}{labiles} et inertes en 1952 et, par
son expérience chrome--cobalt de 1953, a établi le mécanisme de sphère
interne~; il a reçu le prix Nobel de chimie 1983. Rudolph Marcus a établi à
partir de 1956 la théorie du transfert d'électron qui porte son nom, et a
reçu le prix Nobel de chimie 1992.
\end{history}

\section{Exercices}

\begin{exercise}[$\star$]\label{exo:b3:complex-mechanisms:1}
Classer comme \omterm{def:b3:complex-mechanisms:labile}{labile} ou inerte, en justifiant~: $\ce{[Cr(H2O)6]^{3+}}$, $\ce{[Co(NH3)6]^{3+}}$, $\ce{[Cr(H2O)6]^{2+}}$,
$\ce{[Cu(H2O)6]^{2+}}$, $\ce{[Zn(H2O)6]^{2+}}$.
\end{exercise}

\begin{exercise}[$\star$]\label{exo:b3:complex-mechanisms:2}
Montrer que la loi de vitesse D se réduit à $v = k_1[\mathrm{ML_5X}]$ pour $[\mathrm Y]$ élevée et à $v = (k_1k_2/k_{-1})
[\mathrm{ML_5X}][\mathrm Y]/[\mathrm X]$ lorsque $k_{-1}[\mathrm X] \gg k_2[\mathrm Y]$.
\end{exercise}

\begin{exercise}[$\star$]\label{exo:b3:complex-mechanisms:3}
Prévoir les signes de $\Delta V^{\ddagger}$ et de $\Delta^{\ddagger}S^\circ$ pour une substitution D et pour une substitution A.
\end{exercise}

\begin{exercise}[$\star$]\label{exo:b3:complex-mechanisms:4}
Donner les produits de $\ce{[PtCl4]^{2-}}$ avec deux \ce{NH3}, et de $\ce{[Pt(NH3)4]^{2+}}$ avec deux \ce{Cl-}.
\end{exercise}

\begin{exercise}[$\star\star$]\label{exo:b3:complex-mechanisms:5}
Pour une substitution sur un aqua-ion, $K_{\mathrm{os}} = \qty{2.0}{L.mol^{-1}}$ et $k_i = \qty{3.0e4}{s^{-1}}$ (données
de l'exercice). Calculer $k_{\mathrm{obs}}$ pour $[\mathrm Y] = 0.10$ et \qty{1.0}{mol/L}, ainsi que sa limite.
\end{exercise}

\begin{exercise}[$\star\star$]\label{exo:b3:complex-mechanisms:6}
Une constante de vitesse double quand la pression passe de 0.1 à
\qty{100}{MPa} à \qty{298}{K}. Calculer $\Delta V^{\ddagger}$ et conclure.
\end{exercise}

\begin{exercise}[$\star\star$]\label{exo:b3:complex-mechanisms:7}
Dans $\ce{[PtCl3(PEt3)]-}$, la liaison Pt--Cl en trans de \ce{PEt3} mesure \qty{2.42}{\angstrom} et les deux autres \qty{2.32}{\angstrom}
(données de l'exercice). Interpréter, et prévoir quel chlorure est remplacé
en premier.
\end{exercise}

\begin{exercise}[$\star\star$]\label{exo:b3:complex-mechanisms:8}
Énumérer les observations expérimentales qui montreraient qu'un transfert
d'électron est de sphère interne.
\end{exercise}

\begin{exercise}[$\star\star$]\label{exo:b3:complex-mechanisms:9}
Pour $\lambda = \qty{80}{kJ/mol}$, calculer $\Delta G^{\ddagger}$ pour $\Delta G^\circ = 0$, $-20$, $-80$ et \qty{-120}{kJ/mol}.
\end{exercise}

\begin{exercise}[$\star\star\star$]\label{exo:b3:complex-mechanisms:10}
Deux couples ont des constantes d'auto-échange $k_{11} = 4.0$ et $k_{22} = \qty{2.0e5}{L.mol^{-1}.s^{-1}}$, et
leur réaction croisée a $K_{12} = \num{1.0e4}$ (données de l'exercice). Estimer $k_{12}$ ($f = 1$).
\end{exercise}

\begin{exercise}[$\star\star\star$]\label{exo:b3:complex-mechanisms:11}
À l'aide du modèle purement $\sigma$, calculer les énergies d'activation de
champ des ligands des ions $d^5$, $d^7$ et $d^8$ haut spin, et classer
\ce{Mn^{2+}}, \ce{Co^{2+}} et \ce{Ni^{2+}} selon la vitesse d'échange d'eau
attendue.
\end{exercise}

\begin{exercise}[$\star\star\star$]\label{exo:b3:complex-mechanisms:12}
Pourquoi la \omterm{def:b3:complex-mechanisms:inverted}{région inversée} est-elle difficile à observer pour des
réactions entre ions qui diffusent l'un vers l'autre en solution, et
comment des molécules rigides donneur--accepteur l'ont-elles révélée~?
\end{exercise}

\section{Problème~: fabriquer le cisplatine, pas le transplatine}

\begin{problem}[{Problème du week-end --- fabriquer le cisplatine, pas le
transplatine~: l'ion plan carré du platine(II), une synthèse planifiée grâce
à l'effet trans, l'aquation du médicament, et pourquoi il est activé dans
les cellules}]
\label{pb:b3:complex-mechanisms:1}
% ledger: ghs4:K2PtCl4
Le cisplatine, $cis$-$\ce{[PtCl2(NH3)2]}$, se prépare à partir de $\ce{K2[PtCl4]}$. Données du problème~: à
\qty{37}{\celsius}, la première aquation, $\ce{[PtCl2(NH3)2] + H2O -> [PtCl(H2O)(NH3)2]+ + Cl-}$, a la
constante de vitesse $k = \qty{9.0e-5}{s^{-1}}$ et la constante d'équilibre $K = \qty{3.0e-3}{mol/L}$~; la
concentration en chlorure vaut environ \qty{100}{mM} dans le plasma sanguin
et \qty{4}{mM} dans les cellules.

\medskip\noindent\textbf{Partie I --- Le platine(II).}
\begin{enumerate}
  \item Donner le nombre d'électrons $d$ de \ce{Pt^{2+}}.
  \item Pourquoi ses complexes sont-ils plans carrés et diamagnétiques~?
  \item Pourquoi leurs substitutions sont-elles associatives~?
  \item Écrire la loi de vitesse à deux termes et expliquer ses deux termes.
  \item Le médicament est-il \omterm{def:b3:complex-mechanisms:labile}{labile} ou inerte à l'échelle d'une journée~?
\end{enumerate}

\medskip\noindent\textbf{Partie II --- La synthèse.}
\begin{enumerate}[resume]
  \item Que donne $\ce{K2[PtCl4]}$ avec deux \ce{NH3}~? Pourquoi~?
  \item Cette voie directe donne un produit impur. La voie de Dhara
        convertit d'abord $\ce{[PtCl4]^{2-}}$ en $\ce{[PtI4]^{2-}}$ par un excès de KI. Pourquoi
        l'iodure est-il utile~?
  \item L'ajout de \ce{NH3} donne $cis$-$\ce{[PtI2(NH3)2]}$. Expliquer la stéréochimie.
  \item Les iodures sont ensuite éliminés par le nitrate d'argent dans
        l'eau. Que se forme-t-il~?
  \item L'ajout de KCl donne le cisplatine. Pourquoi la configuration ne
        change-t-elle pas~?
  \item Comment préparer plutôt le transplatine~?
  \item Pourquoi le transplatine est-il inutile comme médicament, bien qu'il
        réagisse lui aussi~?
\end{enumerate}

\medskip\noindent\textbf{Partie III --- L'aquation.}
\begin{enumerate}[resume]
  \item Calculer le temps de demi-réaction de la première aquation.
  \item Calculer la fraction aquée au bout de \qty{2.0}{h} dans une solution
        sans chlorure.
  \item Pourquoi l'aquacomplexe est-il la forme réactive vis-à-vis de
        l'ADN~?
  \item L'aquation est-elle dissociative ou associative~? Quel serait son
        \omterm{def:b3:complex-mechanisms:activation-volume}{volume d'activation}~?
  \item Pourquoi une solution injectable du médicament est-elle préparée
        dans du sérum physiologique~?
\end{enumerate}

\medskip\noindent\textbf{Partie IV --- Plasma et cellule.}
\begin{enumerate}[resume]
  \item Écrire la fraction à l'équilibre de l'aquacomplexe en fonction de
        $[\ce{Cl-}]$.
  \item La calculer dans le plasma.
  \item La calculer dans une cellule.
  \item Calculer le rapport.
  \item Pourquoi le médicament agit-il donc surtout dans les cellules~?
  \item Quels autres ligands présents dans une cellule peuvent se lier au
        platine et désactiver le médicament~?
  \item Pourquoi les deux ligands ammine restent-ils liés tout du long~?
  \item Énoncer le résultat~: le rapport de la fraction d'aquacomplexe dans
        une cellule à celle dans le plasma.
\end{enumerate}
\end{problem}
